% =====================================================================
%  cloud_mechanisms.tex
% =====================================================================
\chapter{Cloud mechanisms}
\label{ch:mechanisms}

\section{Infrastructure mechanisms}

\subsection{Logical network perimeter}

The logical network perimeter works at two levels. The outer perimeter
separates the system from the other consumers of the cloud, while the inner
perimeters divide it into the zones described in \cref{sec:zones}.
The data zone is protected by both firewalls and can only be reached from
the application zone. The training zone only receives the pseudonymized export
and publishes new models to the model repository.

\subsection{Virtual server}

Each service runs on its own virtual servers, created from hardened images.
In this way every service is sized independently and a fault in one service
does not affect the others.

\subsection{Cloud storage device}

The data are stored in different devices according to how they are used
and the level of security they require:

\begin{itemize}
    \item \textbf{time series}: stores the measured glucose values, one sample every five minutes per patient, about 167 writes per second;
    \item \textbf{relational database}: stores patient information, doctor assignments, diets,
          meals, forecasts, flags and interventions;
    \item \textbf{model repository}: stores the different model versions, so that every forecast
          can be traced to a model and a previous version can be restored;
    \item \textbf{aggregated view}: contains statistics without personal information and is the only storage reachable by the manager service;
    \item \textbf{pseudonymized copy}: is stored in the training zone and contains the data needed for training without direct identifiers.
\end{itemize}

\subsection{Resource replication}

Resource replication is used both for scaling and for recovery after a failure.
It creates new instances of virtual servers when needed and keeps a synchronized
copy of each storage device in the data zone.

\section{Specialized mechanisms}

\subsection{Load balancer}

The load balancer is the main entry point for all users. It handles HTTPS/TLS
connections, applies a rate limit to each client to help prevent denial-of-service
attacks and distributes requests across the available instances.

It also contacts the IAM to check the user's identity and role before forwarding
a request. Based on this information, the request is sent only to the service
assigned to that role (for example, a patient can never reach the diabetologist or the manager service).

\subsection{Automated scaling listener}


The scaling listener monitors the load of each service and adds or removes
instances when predefined thresholds are reached. The workload changes
throughout the day, with higher demand around breakfast, lunch and dinner.
To avoid wasting resources during quieter periods or having too few instances
during peak hours, the system can adjust the number of instances automatically.

At least two instances of each online service are always kept
available to support failover.

\subsection{Failover system}


The architecture uses an \textit{active-active} configuration, where all instances
handle requests at the same time. If one of them fails, the others continue
serving users while a replacement instance is started. 

We chose this approach
over active-passive because the instances are already needed to handle the
workload. There is also no need to wait for a standby instance to take over,
which helps keep outages short.

Storage devices also have a synchronized
replica ready to replace the primary one if it fails.

\subsection{SLA monitor}

The SLA monitor keeps track of the uptime and response time of each service.
It also sends an alert when a measured value gets close to its allowed limit.
The collected metrics are available to the manager service and displayed
in the center manager's dashboard.

\subsection{Audit monitor}

The audit monitor collects logs of all accesses to health data (who accessed which data, and
when) in a log that the services cannot modify, making it easier to identify
any improper access.

\section{Security mechanisms}
\label{sec:mech-security}

\subsection{Firewalls}


The external firewall allows only HTTPS traffic to reach the load balancer,
and sensor traffic to reach the CGM ingestion component. 

The internal firewall, instead,
controls the connections from the DMZ to the internal zones.

This prevents the
internal components from being directly exposed to the Internet.

\subsection{Encryption}

TLS protects the data in transit against eavesdropping and malicious
intermediaries. Data stored in the devices of the
data zone are also encrypted at rest. 

\subsection{Identity and access management}


The IAM stores user accounts, credentials and roles. It also handles the
identity checks requested by the load balancer. Passwords are always stored only as hashes. 

The load balancer checks the user's role,
while the diabetologist service performs an additional check to make sure
that the requested patient is actually assigned to that doctor.



\subsection{Hardened virtual server images}


Every virtual server starts from an image without unnecessary software,
services and accounts, and with no known critical or high vulnerabilities
(verified through automated vulnerability scanning). Instances created by
scaling or failover have the same protection level.


\subsection{Pseudonymization}

The data exported to the training zone have no direct identifiers. Since the training process does not need to know who the patients
are, this reduces the risk of exposing their identities if the training zone
is compromised.